
\section{Module Organization}
\label{sec:module-organization}
% TODO: Should this section be moved in the Getting Started Tutorial?

In this section we describe the recommended structure of a basic CubedOS module. Modules are not
required to follow this structure exactly; special needs often require special handling.
However, the structure we describe here is appropriate for many modules. The CubedOS source
distribution includes a template module, with comments, that you can use as a starting point for
your own modules.

A module consists of a hierarchy of packages formed by a root package and (at least) two child
packages. The root package can be used to hold module-wide declarations as appropriate. In many
cases the root package may not require a body, although module-wide subprograms might usefully
be declared in the root package.

For a module named |My_Module| we recommend the following two child packages with possible
additional packages as follows:

\begin{itemize}
\item |My_Module.Messages| is a public child package containing the server task that loops
infinitely reading the module mailbox, processing the messages it finds there. This loop drives
the activity of the module, processing exactly one message at a time. No other task in the
module, or indeed in the entire system, should read from the module's mailbox. Furthermore, the
server task should receive a message in only one place in the program. This highly stylized
structure is conceptually simple and easy to analyze.

\item |My_Module.API| is a public child package that contains helper subprograms for
constructing and parsing module messages. These helper subprograms only manipulate message
records. They do no actual message processing. Other modules will call these subprograms to
build messages sent to and parse messages received from |My_Module|. The |API| package may also
include type definitions related to the module's public interface.

A separate tool, Mercury, can be used to autogenerate the |API| package from an interface
definition. Mercury is described in Chapter~\ref{chapt:mercury}.

\item Optionally, |My_Module.Internals| is a private child package that contains the bulk of the
module's processing logic. In simple cases, this package may not be necessary and all processing
of messages can be in the |Messages| package. However, complex modules may find it useful to
break out the bulk of the processing logic into an |Internals| package to separate that logic
from the CubedOS infrastructure associated with message handling. The code in this package is
not accessible from outside the module (since it is a private child package), and it can assume
it is called by only one task unless additional, internal tasks are created by the module
developer.

In even more complex cases, it may make sense for there to be multiple private "internals"
packages, or even a private internal package hierarchy. Organizing the internal logic of a
module is ultimately at the discretion of the module developer.
\end{itemize}

Although this structure seems complex it has several advantages.

The bulk of the module's logic is in packages (both |API| and |Internals|) that, in most cases,
have no tasking structures. Typically, for modules with no internal tasks, all the code in any
|Internals| package is executed by only one task, simplifying development. The |API| encoding
and decoding subprograms might be called by multiple tasks, but they are relatively simple, with
no global effects, and thus easy to make task-safe.

Because the |Internals| package is a private child package, its resources cannot be directly
accessed by the rest of the program.

Using the module API is highly stylized and entails calling an API function to build a message
that is then sent to the module's mailbox. This is most conveniently done using |Route_Message|
in the |Message_Manager| package:

\begin{lstlisting}
   Route_Message(My_Module.API.Operation_Request_Encode(1, 2, 3));
\end{lstlisting}

Here |Operation_Request_Encode| is a function that creates a message causing |My_Module| to
perform some designated operation. The name of the function should be reflective of the
operation being done. We recommend using the convention of adding ``Request'' as a suffix to the
names of messages making requests sent to the module, and adding ``Reply'' as a suffix to the
names of messages that are replies from the module. This aligns with the concept that modules
are, in effect, servers that accept requests and return replies.

The arguments to the function (1, 2, and 3 in this case) are arguments to the operation and are
packed into a message following XDR rules. The Mercury tool auto-generates |API| bodies that
take care of this. Note that |Route_Message| handles looking up a module's ID in the name
resolver package. In the case of a multi-domain application, |Route_Message| will automatically
direct the message to some underlying message router, as needed.

Finally, the architecture of modules is compatible with the Jorvik restrictions. In particular,
a module's server task is a single, library level task that runs infinitely, and all
communication with a module is by way of a protected mailbox object with only one entry for
receiving messages.

